\section{A bounded Fourier--Mellin bridge on a specified data line}
\label{sec:frameworkbridge}
The relation between Fourier and Mellin presentations is classical.
The following calculation isolates its noise norm when both signs of
retarded time are available. It is a kinematic transform identity, not
a construction of a holographic dual or an entanglement formula.

Let $f\in C_c^\infty((0,\infty);\mathbb C)$ and write
\[
 F(q)=\int_0^\infty\omega^{q-1}f(\omega)\,d\omega,
 \qquad
 C_\pm(u)=\int_0^\infty e^{\mp i\omega u}f(\omega)\,d\omega,
 \quad u>0.
\]
Thus $C_+(u)$ and $C_-(u)$ are the restrictions of the same
positive-energy Fourier transform to positive and negative retarded
time, respectively. Having both restrictions is a data assumption.
Define $U_\pm(s)=\int_0^\infty u^{s-1}C_\pm(u)\,du$.
For $0<\Re s<1$ these integrals converge absolutely: $C_\pm$ are
bounded near zero and decay faster than every inverse power at infinity.

\begin{proposition}[Two-channel transform and noise norm]
\label{prop:frameworkbridge}
For $0<\Re s<1$,
\begin{equation}\label{eq:frameworkbridge}
 U_\pm(s)=k_\pm(s)F(1-s),\qquad
 k_\pm(s)=\Gamma(s)e^{\mp i\pi s/2}.
\end{equation}
On $s=\frac12+i\nu$, $\nu\in\mathbb R$, one has
\begin{equation}\label{eq:jointkernel}
 |k_+(s)|^2+|k_-(s)|^2=2\pi.
\end{equation}
Consequently, observations $V_\pm=U_\pm+e_\pm$ admit the estimator
\begin{equation}\label{eq:frameworkestimator}
 \widehat F(1-s)=\frac{\overline{k_+(s)}V_+(s)
                         +\overline{k_-(s)}V_-(s)}{2\pi}
\end{equation}
with the pointwise sharp linear noise bound
\begin{equation}\label{eq:frameworknoise}
 |\widehat F(1-s)-F(1-s)|
 \le\frac{\sqrt{|e_+(s)|^2+|e_-(s)|^2}}{\sqrt{2\pi}}.
\end{equation}
The corresponding operator norm on joint $L^2(d\nu)$ errors is also
$1/\sqrt{2\pi}$. Moreover,
\[
 \|U_+\|_{L^2(d\nu)}^2+\|U_-\|_{L^2(d\nu)}^2
 =2\pi\|F(\tfrac12-i\nu)\|_{L^2(d\nu)}^2.
\]
\end{proposition}
\begin{proof}
Insert an Abel factor $e^{-\varepsilon u}$, $\varepsilon>0$.
Absolute integrability permits Fubini and the elementary gamma integral:
\[
 \int_0^\infty u^{s-1}e^{-\varepsilon u}C_\pm(u)\,du
 =\Gamma(s)\int_0^\infty
       f(\omega)(\varepsilon\pm i\omega)^{-s}\,d\omega.
\]
As $\varepsilon\downarrow0$, dominated convergence applies on the
left using $u^{\Re s-1}|C_\pm(u)|$, and on the right because the
compact support stays away from zero. Principal logarithms give
$(\pm i\omega)^{-s}=\omega^{-s}e^{\mp i\pi s/2}$, proving
\eqref{eq:frameworkbridge}. The reflection identity yields
\[
 |\Gamma(\tfrac12+i\nu)|^2=\frac\pi{\cosh(\pi\nu)}.
\]
Multiplication by $e^{\pi\nu}+e^{-\pi\nu}$ proves
\eqref{eq:jointkernel}. Equation~\eqref{eq:frameworkestimator}
recovers the exact target in the absence of noise, and Cauchy--Schwarz
proves \eqref{eq:frameworknoise}. Equality holds when the two-component
error is aligned with $(k_+,k_-)$, so the displayed linear noise norm
is sharp. Integrating proves the $L^2$ statements; the target is in
$L^2$ by the Fourier representation in logarithmic energy.
\end{proof}

\begin{proposition}[One time channel has no uniform inverse bound]
\label{prop:frameworkonechannel}
On the same line,
\[
 |k_+(\tfrac12+i\nu)|^{-2}
   =\frac{1+e^{-2\pi\nu}}{2\pi},\qquad
 |k_-(\tfrac12+i\nu)|^{-2}
   =\frac{1+e^{2\pi\nu}}{2\pi}.
\]
Thus one channel has exponentially large frequencywise inverse
amplification in one direction. No uniform $L^2$ inverse estimate
holds on the stated compact smooth source class using only that channel.
\end{proposition}
\begin{proof}
The formulas follow directly from the preceding gamma identity.
For an actual source-class sequence, choose nonzero
$f_0\in C_c^\infty((0,\infty))$ and
$f_N(\omega)=\omega^{-iN}f_0(\omega)$.
Then $F_N(\tfrac12-i\nu)=F_0(\tfrac12-i(\nu+N))$,
so the target $L^2$ norm is independent of $N$.
After setting $\mu=\nu+N$, the positive-time data norm squared is
\[
 \int_{\mathbb R}|k_+(\tfrac12+i(\mu-N))|^2
                 |F_0(\tfrac12-i\mu)|^2\,d\mu\longrightarrow0.
\]
Indeed, the kernel tends pointwise to zero and is bounded by $2\pi$;
dominated convergence applies. The other channel is treated by the
opposite modulation.
\end{proof}

The line in this calculation is $\Re q=1/2$ for the unweighted
transform $F$. It is not automatically the physical four-dimensional
principal-series line $\Re\Delta=1$. Translating conventions requires
specifying an energy weight and the corresponding observation norm.
The two data channels do not supply a boundary action, operator products,
Ward identity, bulk reconstruction, or an information-theoretic observable.
The result only tests which part of changing transform presentation
is stable under a clearly stated two-channel noise model.
